% Slides — Capítulo 17: Ventos Regionais
\documentclass[aspectratio=169,11pt]{beamer}
\usepackage{../comum/temameteo}
\graphicspath{{../figuras/cap17/}}

\title[Cap. 17 --- Ventos Regionais]{Capítulo 17 --- Ventos Regionais}
\subtitle{Meteorologia Prática}
\author{}
\date{}

\begin{document}

\begin{frame}
  \titlepage
  \vfill
  \atribuicao
\end{frame}

\begin{frame}{Roteiro}
  \tableofcontents
\end{frame}

% ---------------------------------------------------------------
\section{Circulações térmicas}

\begin{frame}{Brisa marítima: a frente que entra do mar}
  \begin{columns}
    \column{0.34\textwidth}
    \begin{itemize}
      \item Terra quente de dia (Bowen, Cap.~3): circuito de
            Bjerknes (T1).
      \item Mini-frente fria avançando 10--50\,km terra adentro —
            gatilho da convecção da tarde (Cap.~14).
      \item À noite, o terral fraco.
    \end{itemize}
    \column{0.33\textwidth}
    \figlivroslide{fig_17_13.jpg}{0.42\textheight}{Fig.~17.13 --- a
      frente de brisa em corte}
    \column{0.33\textwidth}
    \figlivroslide{fig_17_14.jpg}{0.42\textheight}{Fig.~17.14 --- o
      relógio da brisa: Coriolis gira o vetor}
  \end{columns}
\end{frame}

\begin{frame}{Vale e encosta: o dia e a noite da montanha}
  \begin{columns}
    \column{0.4\textwidth}
    \begin{itemize}
      \item Dia: anabático sobe a encosta $+$ vento de vale sobe o
            eixo (nuvens de crista ao meio-dia).
      \item Noite: catabático $+$ vento de montanha descem —
            $u_{eq} = \sqrt{gh\Delta\theta\sin\alpha/C_D\theta}$
            (T2).
      \item Piscinas frias de fundo de vale: as geadas do café
            (Cap.~18).
    \end{itemize}
    \column{0.6\textwidth}
    \figlivroslide{fig_17_11.jpg}{0.52\textheight}{Fig.~17.11 ---
      anabáticos e ventos de vale diurnos}
  \end{columns}
\end{frame}

% ---------------------------------------------------------------
\section{Ondas de montanha}

\begin{frame}{Os regimes de Froude num desenho}
  \begin{columns}
    \column{0.34\textwidth}
    \begin{itemize}
      \item $Fr \ll 1$: contorna/bloqueia; $Fr \sim 1$: ondas e
            rotores; $Fr \gg 1$: cavidade turbulenta.
      \item A serra ``toca'' a mola $N$ do Cap.~5.
      \item Rotores sob as cristas: perigo de aviação.
    \end{itemize}
    \column{0.33\textwidth}
    \figlivroslide{fig_17_30e.jpg}{0.4\textheight}{Fig.~17.30 ---
      $Fr \sim 1$: ondas e rotores a sotavento}
    \column{0.33\textwidth}
    \figlivroslide{fig_17_30g.jpg}{0.4\textheight}{Fig.~17.30 ---
      turbulência na esteira}
  \end{columns}
\end{frame}

\begin{frame}{Lenticulares: as cristas fotografáveis}
  \begin{columns}
    \column{0.4\textwidth}
    \eqdestaque{$\lambda_z = \dfrac{2\pi U}{N} \sim 7$\,km}
    \vspace{0.4em}
    \begin{itemize}
      \item Nuvem PARADA com vento forte através: condensa na crista,
            evapora no vale (Cap.~6).
      \item Cristas inclinam contra o vento com a altura: energia
            propagando para cima (T3; Caso~3).
      \item Quebra em altitude: CAT (Cap.~5) e arrasto de onda
            (Cap.~20).
    \end{itemize}
    \column{0.6\textwidth}
    \figlivroslide{fig_17_30c.jpg}{0.5\textheight}{Fig.~17.30 --- LCL
      cruzado nas cristas: o trem de lenticulares}
  \end{columns}
\end{frame}

% ---------------------------------------------------------------
\section{Hidráulica de sotavento}

\begin{frame}{O salto hidráulico da serra}
  \begin{columns}
    \column{0.4\textwidth}
    \eqdestaque{$Fr = \dfrac{U}{\sqrt{g'h}}$}
    \vspace{0.4em}
    \begin{itemize}
      \item Transição crítica no topo $\Rightarrow$ a camada
            despenca supercrítica e fecha num salto turbulento (T4;
            Caso~4).
      \item A física da pia da cozinha — e da Bora e dos windstorms
            de Boulder.
      \item Gap winds: Bernoulli no desfiladeiro.
    \end{itemize}
    \column{0.6\textwidth}
    \figlivroslide{fig_17_18.jpg}{0.52\textheight}{Fig.~17.18 ---
      sub$\to$supercrítico$\to$salto: o vendaval de sotavento}
  \end{columns}
\end{frame}

\begin{frame}{Föhn: desce quente e seco}
  \begin{columns}
    \column{0.34\textwidth}
    \begin{itemize}
      \item Sobe úmido (chove na serra), desce TUDO pelo
            $\Gamma_d$: $+\Delta T$ e ar seco (T5).
      \item O $L_v$ que ficou a barlavento paga a conta (Caps.~3--4).
      \item A parede de föhn no cume e o céu limpo a sotavento.
    \end{itemize}
    \column{0.33\textwidth}
    \figlivroslide{fig_17_38.jpg}{0.42\textheight}{Fig.~17.38 --- a
      parede de föhn}
    \column{0.33\textwidth}
    \figlivroslide{fig_17_39.jpg}{0.42\textheight}{Fig.~17.39 --- o
      caminho do föhn no emagrama}
  \end{columns}
\end{frame}

% ---------------------------------------------------------------
\section{Vento como recurso e a cidade}

\begin{frame}{Energia eólica: o cubo do vento}
  \begin{columns}
    \column{0.34\textwidth}
    \begin{itemize}
      \item Potência $\propto \rho A v^3$; limite de Betz:
            $E_o^{max} = 16/27 \approx 59\%$.
      \item Perfil log (Cap.~18) decide a altura da torre.
      \item Curva real: cut-in, rated, cut-out.
    \end{itemize}
    \column{0.33\textwidth}
    \figlivroslide{fig_17_5a.jpg}{0.42\textheight}{Fig.~17.5 ---
      eficiência vs.\ razão de velocidades (Betz)}
    \column{0.33\textwidth}
    \figlivroslide{fig_17_5b.jpg}{0.42\textheight}{Fig.~17.5 --- curva
      de potência de uma turbina}
  \end{columns}
\end{frame}

\begin{frame}{A ilha de calor urbana}
  \begin{columns}
    \column{0.4\textwidth}
    \begin{itemize}
      \item Concreto guarda calor (Cap.~3), cânions prendem IV
            (Cap.~2), impermeabilização mata o $F_E$: cidade
            2--6\,K mais quente à noite.
      \item $\Delta T_{max}$ cresce com a razão altura/largura das
            ruas.
      \item A pluma urbana e a brisa rural fecham o quadro
            (Cap.~19).
    \end{itemize}
    \column{0.6\textwidth}
    \figlivroslide{fig_17_41c.jpg}{0.48\textheight}{Fig.~17.41 --- a
      ilha de calor em planta}
  \end{columns}
\end{frame}

% ---------------------------------------------------------------
\section{Síntese e laboratório}

\begin{frame}{O mapa do capítulo}
  \eqdestaque{$\oint dP/\rho$ (brisas) \quad
    $u_{eq}=\sqrt{\tfrac{gh\Delta\theta\sin\alpha}{C_D\theta}}$ \quad
    $\lambda_z = \tfrac{2\pi U}{N}$ \quad
    $Fr = \tfrac{U}{\sqrt{g'h}}$}
  \vspace{0.6em}
  Térmicos de dia, catabáticos de noite, ondas e saltos onde há serra —
  e o vento como recurso ($v^3$!) fechando a conta.
\end{frame}

\begin{frame}{Laboratório numérico --- notebook do Cap.~17}
  \texttt{notebooks/cap17\_ventosregionais.ipynb}
  \begin{enumerate}
    \item Brisa com Coriolis: EDO $+$ hodógrafo diurno.
    \item Catabático: da serra à Antártida.
    \item Ondas de montanha lineares (morro-sino).
    \item Água rasa sobre a serra: regimes de Froude.
  \end{enumerate}
  \vspace{0.4em}
  \textbf{Exercícios}: T1--T5 e N1--N4 nas notas; células reservadas no
  notebook.
  \vfill
  \atribuicao
\end{frame}

\end{document}
